\chapter*{Introduction Générale}
\addcontentsline{toc}{chapter}{Introduction Générale}

\section*{Contexte et Motivation}
À l'ère de la transformation numérique, les outils technologiques redéfinissent en profondeur la majorité des secteurs d'activité humaine. Le domaine du sport, de la santé et du bien-être, historiquement ancré dans des pratiques présentielles et manuelles, connaît une véritable révolution digitale. Les professionnels de la préparation physique, les nutritionnistes et les coachs personnels sont confrontés à une clientèle de plus en plus exigeante, qui réclame un suivi personnalisé, des retours en temps réel, et une approche scientifique basée sur la donnée pour atteindre ses objectifs.

Cependant, la gestion de cette relation "Coach-Athlète" souffre d'un manque criant de standardisation technologique. Aujourd'hui, un professionnel du sport s'appuie généralement sur un empilement chaotique d'outils génériques : des tableurs complexes pour la planification des macro-cycles d'entraînement, des applications tierces pour le suivi calorique, et des messageries grand public pour la communication quotidienne. Cette fragmentation entraîne inévitablement des pertes d'informations, une augmentation drastique de la charge mentale et administrative pour le coach, et une détérioration de l'expérience utilisateur pour l'athlète, qui se retrouve livré à lui-même lors de ses séances en salle.

Face à ces constats, notre projet de fin d'études se donne pour ambition de répondre à une question centrale : \textbf{Comment concevoir, modéliser et développer une solution logicielle unifiée, moderne et évolutive, capable de digitaliser l'intégralité du cycle de vie du coaching sportif professionnel ?}

\section*{Objectifs du Projet}
Pour répondre à cette problématique, nous avons conceptualisé "CoachingApp", un écosystème logiciel composé de trois éléments fondamentaux :
\begin{enumerate}
    \item \textbf{Une application mobile (Front-end Athlète/Coach) :} Pensée pour la mobilité, elle permet à l'athlète de suivre ses séances de manière interactive grâce à un lecteur d'entraînement chronométré (Workout Player), de logger son alimentation, et de communiquer avec son mentor.
    \item \textbf{Un tableau de bord web (Front-end Admin) :} Dédié à la gestion de la plateforme, à la création de catalogues d'exercices standardisés et à l'analyse des métriques d'usage.
    \item \textbf{Un serveur centralisé (Back-end) :} Le cœur du système, exposant une API RESTful sécurisée, gérant les bases de données relationnelles complexes, et permettant une communication en temps réel et des suggestions basées sur l'Intelligence Artificielle.
\end{enumerate}

\section*{Plan du Rapport}
Le présent rapport de projet de fin d'études détaille la démarche scientifique, technique et méthodologique adoptée pour mener à bien cette mission. Il est structuré en six chapitres cohérents, suivant la progression logique du cycle de vie du développement logiciel :

\textbf{Le premier chapitre} dressera le contexte général du projet. Nous y exposerons la problématique détaillée, une étude comparative rigoureuse des solutions concurrentes existantes, ainsi que le cadre méthodologique (Agile Scrum) qui a régi notre gestion de projet.

\textbf{Le deuxième chapitre} sera consacré à l'analyse et à la spécification des besoins. Nous y définirons les acteurs du système à travers un dictionnaire de cas d'utilisation approfondi. Nous justifierons également, de manière technique, les choix architecturaux et technologiques retenus (Flutter, Angular, Node.js, PostgreSQL).

\textbf{Le troisième chapitre} inaugurera la phase de réalisation en décrivant les accomplissements du Sprint 1. Il couvrira la mise en place de la base de données relationnelle, des mécanismes de sécurité et de l'authentification (JWT), ainsi que le processus de mise en relation entre coachs et athlètes.

\textbf{Le quatrième chapitre} abordera le Sprint 2, qui constitue le cœur métier de l'application. Nous y détaillerons les algorithmes complexes régissant la création des programmes d'entraînement hiérarchisés, le fonctionnement asynchrone du lecteur de séance mobile (Workout Player), ainsi que le module de calcul nutritionnel.

\textbf{Le cinquième chapitre} exposera la réalisation du Sprint 3. Il se concentrera sur l'intégration des fonctionnalités de communication en temps réel (WebSockets), le développement du tableau de bord administrateur (Angular), et l'intégration innovante d'une API d'Intelligence Artificielle (Gemini) comme assistant virtuel.

\textbf{Enfin, le sixième chapitre} détaillera la phase cruciale de l'ingénierie logicielle : la stratégie d'Assurance Qualité (Tests unitaires et d'intégration), le renforcement de la sécurité de l'API (protection OWASP), et les procédures de déploiement (DevOps) vers l'environnement de production en nuage (Cloud).

Ce rapport s'achèvera par une conclusion générale qui dressera le bilan de nos travaux, analysera les défis surmontés et proposera des perspectives d'évolution futures pour la plateforme CoachingApp.
